% SOURCE: Chinese PDF viewer pp.177--184 (top), source Section 2 and six subdivisions.
\section{Relativistic Invariance}
\subsection*{1) Lorentz invariance}
Consider two coordinate systems $O$ and $O'$, in which the Dirac equations are respectively
\[
\ii\gamma^\mu\frac{\partial}{\partial x^\mu}\psi(x)-m\psi(x)=0,
\qquad
\ii\gamma^\mu\frac{\partial}{\partial x'^\mu}\psi'(x')-m\psi'(x')=0.
\]
Assume that the coordinates $x^\mu$ and wavefunction $\psi$ are related in the two coordinate systems by
\[
x'=\Lambda x,\qquad \psi'(x')=S(\Lambda)\psi(x).
\]
Here $\Lambda$ and $S(\Lambda)$ are $4\times4$ matrices and are not functions of the coordinates. Since Lorentz scalars remain invariant under a transformation,
\[
x'_\mu x'^\mu=x_\mu x^\mu,
\]
the matrix $\Lambda$ satisfies
\[
\Lambda^\nu{}_{\mu}\Lambda^\mu{}_{\rho}=g^\nu{}_{\rho}
\]
(in metric-contracted index notation). Taking determinants of the Lorentz-invariance condition gives
\[
(\det\Lambda)^2=1,\qquad\det\Lambda=\pm1.
\]
A transformation with $\det\Lambda=+1$ is called a proper Lorentz transformation; one with $\det\Lambda=-1$ is called an improper Lorentz transformation.

Under the above transformation, the Dirac equation in the $O'$ system is
\begin{align*}
\ii\gamma^\mu\frac{\partial x^\nu}{\partial x'^\mu}
\frac{\partial}{\partial x^\nu}S\psi(x)-mS\psi(x)&=0,\\
\ii\gamma^\mu(\Lambda^{-1})^\nu{}_{\mu}
\frac{\partial}{\partial x^\nu}S\psi(x)-mS\psi(x)&=0.
\end{align*}
Multiplying this equation from the left by $S^{-1}$ gives
\[
\ii S^{-1}\gamma^\mu S(\Lambda^{-1})^\nu{}_{\mu}
\frac{\partial\psi(x)}{\partial x^\nu}-m\psi(x)=0.
\]
For the Dirac equation to be Lorentz invariant, i.e. for the $O'$ and $O$ equations to be equivalent, the transformation $S$ must satisfy
\[
S\gamma^\mu S^{-1}=(\Lambda^{-1})^\mu{}_{\nu}\gamma^\nu,
\qquad\text{or}\qquad S^{-1}\gamma^\mu S=\Lambda^\mu{}_{\nu}\gamma^\nu.
\]
This is the restriction on the transformation $S$ imposed by relativistic invariance of the Dirac equation.

For example, consider an infinitesimal transformation
\[
\Lambda^\mu{}_{\nu}=g^\mu{}_{\nu}+\omega^\mu{}_{\nu},
\qquad \omega^{\mu\nu}=-\omega^{\nu\mu},
\]
where $\omega^{\mu\nu}$ is completely antisymmetric in its two indices, and
\[
(\Lambda^{-1})^\mu{}_{\nu}=g^\mu{}_{\nu}-\omega^\mu{}_{\nu},
\qquad\omega^\mu{}_{\nu}=\omega^{\mu\sigma}g_{\sigma\nu}.
\]
Take
\[
S=I-\frac{\ii}{4}\sigma_{\mu\nu}\omega^{\mu\nu},\qquad
S^{-1}=I+\frac{\ii}{4}\sigma_{\mu\nu}\omega^{\mu\nu},
\]
where the $\sigma_{\mu\nu}$ are $4\times4$ matrices to be determined. Substituting into the condition on $S$, one obtains
\begin{align*}
\left(I-\frac{\ii}{4}\sigma_{\mu\nu}\omega^{\mu\nu}\right)
\gamma^\rho
\left(I+\frac{\ii}{4}\sigma_{\mu\nu}\omega^{\mu\nu}\right)
&=\gamma^\nu(g^\rho{}_{\nu}-\omega^\rho{}_{\nu}),\\
\frac{\ii}{4}(\gamma^\rho\sigma_{\mu\nu}-\sigma_{\mu\nu}\gamma^\rho)\omega^{\mu\nu}
&=-\gamma_\nu\omega^{\rho\nu},\\
\frac{\ii}{4}[\gamma^\rho,\sigma_{\mu\nu}]\omega^{\mu\nu}
&=\frac12(-\gamma_\nu\omega^{\rho\nu}-\gamma_\mu\omega^{\mu\rho})\\
&=-\frac12(\gamma_\nu g^\rho{}_{\mu}-\gamma_\mu g^\rho{}_{\nu})\omega^{\mu\nu}.
\end{align*}
Hence $\sigma_{\mu\nu}$ can be chosen as
\[
\sigma_{\mu\nu}=\frac{\ii}{2}[\gamma_\mu,\gamma_\nu].
\]
Clearly $\sigma_{\mu\nu}=-\sigma_{\nu\mu}$ is an antisymmetric tensor matrix.

\subsection*{2) Spin}
Particle spin is automatically contained in the Dirac equation.

For an infinitesimal Lorentz transformation,
\[
\psi'(x')=\left(I-\frac{\ii}{4}\sigma_{\mu\nu}\omega^{\mu\nu}\right)\psi(x)
=\left(I-\frac{\ii}{4}\sigma_{\mu\nu}\omega^{\mu\nu}\right)\psi(\Lambda^{-1}x').
\]
Performing a Taylor expansion,
\begin{align*}
\psi'(x')
&=\left(I-\frac{\ii}{4}\sigma_{\mu\nu}\omega^{\mu\nu}\right)
\left(1+x'_\mu\omega^{\mu\nu}\partial'_\nu\right)\psi(x'),\\
\psi'(x)
&=\left(I-\frac{\ii}{4}\sigma_{\mu\nu}\omega^{\mu\nu}\right)
\left(1+x_\mu\omega^{\mu\nu}\partial_\nu\right)\psi(x)\\
&=\left(I-\frac{\ii}{4}\sigma_{\mu\nu}\omega^{\mu\nu}
+x_\mu\omega^{\mu\nu}\partial_\nu\right)\psi(x)\\
&=\left[I-\frac{\ii}{4}\sigma_{\mu\nu}\omega^{\mu\nu}
+\frac12(x_\mu\partial_\nu-x_\nu\partial_\mu)\omega^{\mu\nu}\right]\psi(x).
\end{align*}
Terms of second order in the infinitesimal parameters have been neglected.

For an infinitesimal rotation,
\[
\psi'(x)=\left(I-\frac{\ii}{2}J_{\mu\nu}\omega^{\mu\nu}\right)\psi(x),
\]
where the angular-momentum tensor $J_{\mu\nu}$ generates rotations. Comparison with the preceding equation gives
\[
J_{\mu\nu}=\frac12\sigma_{\mu\nu}
+\ii(x_\mu\partial_\nu-x_\nu\partial_\mu).
\]
Writing the angular-momentum components as
\[
J_1=J_{23},\qquad J_2=J_{31},\qquad J_3=J_{12},
\]
we obtain the total angular momentum
\[
\vect J=\frac12\begin{pmatrix}\vect\sigma&0\\0&\vect\sigma\end{pmatrix}
+\vect r\times\vect p
=\frac12\vect\Sigma+\vect L,
\qquad \vect p=-\ii\vect\nabla.
\]
Thus the total angular momentum $\vect J$ consists of orbital angular momentum $\vect L$ and spin angular momentum $\vect S=\frac12\vect\Sigma$.

The fact that the Dirac equation automatically contains spin is a relativistic effect. Specifically, this is a consequence of the Dirac equation containing first-order derivatives; the Klein--Gordon equation is second order and does not contain spin.

\subsection*{3) Spatial-inversion invariance ($P$ transformation)}
As shown above, invariance of the Dirac equation under $\psi'(x')=S\psi(x)$ requires
\[
S^{-1}\gamma^\mu S=\Lambda^\mu{}_{\nu}\gamma^\nu.
\]
For continuous transformations we can take $S=I-\frac{\ii}{4}\sigma_{\mu\nu}\omega^{\mu\nu}$, where $\sigma_{\mu\nu}=\frac{\ii}{2}[\gamma_\mu,\gamma_\nu]$. For a discontinuous transformation, such as parity,
\[
x'=\Lambda x,\qquad
\Lambda^{\mu\nu}=\operatorname{diag}(1,-1,-1,-1)=g^{\mu\nu}.
\]
Choose
\[
S=P=\gamma^0.
\]
Then
\[
S^{-1}\gamma^\mu S=\gamma^0\gamma^\mu\gamma^0
=g^{\mu\mu}\gamma^\mu=\Lambda^\mu{}_{\nu}\gamma^\nu
\quad\text{(no sum in $\mu$)}.
\]
This satisfies the condition on $S$, showing that the Dirac equation is invariant under spatial inversion.

\subsection*{4) Time-reversal invariance ($T$ transformation)}
For a time-reversal transformation,
\[
x'=\Lambda x,\qquad
\Lambda^{\mu\nu}=\operatorname{diag}(-1,1,1,1)=-g^{\mu\nu}.
\]
Choose
\[
S=T=\gamma_1\gamma_2\gamma_3.
\]
Since
\begin{align*}
SS&=\gamma_1\gamma_2\gamma_3\gamma_1\gamma_2\gamma_3\\
&=-\gamma_1\gamma_2\gamma_1\gamma_3\gamma_2\gamma_3
=\gamma_1\gamma_1\gamma_2\gamma_3\gamma_2\gamma_3\\
&=-\gamma_1\gamma_1\gamma_2\gamma_2\gamma_3\gamma_3=I,
\end{align*}
we have $S^{-1}=S$, and
\[
S^{-1}\gamma^\mu S=S\gamma^\mu S
=S\gamma^\mu\gamma_1\gamma_2\gamma_3
=\Lambda^\mu{}_{\nu}S\gamma_1\gamma_2\gamma_3\gamma^\nu
=\Lambda^\mu{}_{\nu}SS\gamma^\nu
=\Lambda^\mu{}_{\nu}\gamma^\nu.
\]
This satisfies the condition on $S$ and shows that the Dirac equation is time-reversal invariant.

\subsection*{5) Charge-conjugation invariance ($C$ transformation)}
Introduce an electromagnetic interaction by minimal coupling. Replace the derivative in the Dirac equation by the covariant derivative
\[
\partial_\mu\longrightarrow D_\mu=\partial_\mu+\ii e A_\mu,
\]
where the electromagnetic potential $A_\mu=(A_0,\vect A)$ is an external field, not a wavefunction. The Dirac equation is
\[
(\ii\gamma^\mu\partial_\mu-e\gamma^\mu A_\mu-m)\psi=0.
\]
It remains invariant under a local gauge transformation
\[
\psi(x)\longrightarrow e^{\ii e\alpha(x)}\psi(x),\qquad
A_\mu(x)\longrightarrow A_\mu(x)-\partial_\mu\alpha(x).
\]
Take the Hermitian conjugate of the Dirac equation and multiply from the right by $\gamma^0$:
\[
\overline\psi\left(\ii\gamma^\mu\overleftarrow\partial_\mu
+e\gamma^\mu A_\mu+m\right)=0.
\]
Transpose:
\[
\left(\ii\gamma^{\mu\mathsf T}\partial_\mu
+e\gamma^{\mu\mathsf T}A_\mu+m\right)\overline\psi^{\mathsf T}=0.
\]
Multiply on the left by the transformation matrix $C$:
\[
\left(\ii C\gamma^{\mu\mathsf T}C^{-1}\partial_\mu
+eC\gamma^{\mu\mathsf T}C^{-1}A_\mu+m\right)C\overline\psi^{\mathsf T}=0.
\]
If $C$ satisfies
\[
C\gamma^{\mu\mathsf T}C^{-1}=-\gamma^\mu,
\]
and we define
\[
\psi^c=C\overline\psi^{\mathsf T},
\]
then
\[
(\ii\gamma^\mu\partial_\mu+e\gamma^\mu A_\mu-m)\psi^c=0.
\]
Compared with the Dirac equation for a particle with charge $e$, the wavefunction $\psi^c$ describes a particle of charge $-e$. The transformation $C$ is called charge conjugation, or the particle--antiparticle transformation.

Note that $C$ is an internal transformation, not a spacetime transformation.

\subsection*{6) Bilinear covariants}
Physically observable quantities should be constructed from the Dirac spinor
\[
\psi(x)=\begin{pmatrix}\psi_1(x)\\\psi_2(x)\\\psi_3(x)\\\psi_4(x)\end{pmatrix},
\qquad \overline\psi(x)=\psi^\dagger(x)\gamma^0,
\]
in the form of Dirac bilinears $\overline\psi\Gamma\psi$, with $\Gamma$ a $4\times4$ matrix. What are their Lorentz-transformation properties?

First consider the transformation of $\overline\psi(x)$. Since
\[
\psi'(x')=S\psi(x),
\]
we have
\[
\psi'^\dagger(x')=\psi^\dagger(x)S^\dagger,
\]
and multiplication on the right by $\gamma^0$ gives
\[
\overline\psi'(x')=\overline\psi(x)\gamma^0 S^\dagger\gamma^0.
\]
Since
\[
\gamma^0\sigma_{\mu\nu}^\dagger\gamma^0=\sigma_{\mu\nu},
\qquad
\gamma^0S^\dagger\gamma^0=
\gamma^0\left(I+\frac{\ii}{4}\sigma_{\mu\nu}^\dagger\omega^{\mu\nu}\right)\gamma^0
=I+\frac{\ii}{4}\sigma_{\mu\nu}\omega^{\mu\nu}=S^{-1},
\]
it follows that
\[
\overline\psi'(x')=\overline\psi(x)S^{-1}.
\]
There are 16 linearly independent $4\times4$ matrices, which can be taken as
\[
\Gamma\in\left\{\begin{array}{ll}
I&(1),\\
\gamma_\mu&(4),\\
\sigma_{\mu\nu}&(6),\\
\gamma_5&(1),\\
\gamma_5\gamma_\mu&(4),
\end{array}\right.
\]
where
\[
\gamma_5=\ii\gamma_0\gamma_1\gamma_2\gamma_3.
\]
Its properties include
\[
\{\gamma_5,\gamma_\mu\}=0,\qquad
\gamma_5^2=I,\qquad S^{-1}\gamma_5S=\gamma_5.
\]
Using the transformation properties of $\Gamma$,
\begin{align*}
S^{-1}IS&=I,\\
S^{-1}\gamma_\mu S&=\Lambda_\mu{}^\nu\gamma_\nu,\\
S^{-1}\sigma_{\mu\nu}S&=\Lambda_\mu{}^\rho\Lambda_\nu{}^\sigma\sigma_{\rho\sigma},\\
S^{-1}\gamma_5S&=-\gamma_5,\\
S^{-1}\gamma_5\gamma_\mu S&=-\Lambda_\mu{}^\nu\gamma_5\gamma_\nu,
\end{align*}
we see that the 16 independent Dirac bilinears transform as
\begingroup\small
\begin{align*}
\overline\psi'(x')\psi'(x')&=\overline\psi(x)S^{-1}S\psi(x)
=\overline\psi(x)\psi(x),&&\text{Lorentz scalar},\\
\overline\psi'(x')\gamma^\mu\psi'(x')
&=\overline\psi(x)S^{-1}\gamma^\mu S\psi(x)
=\Lambda^\mu{}_{\nu}\overline\psi(x)\gamma^\nu\psi(x),&&\text{Lorentz vector},\\
\overline\psi'(x')\gamma^5\psi'(x')
&=\overline\psi(x)S^{-1}\gamma^5S\psi(x)
=\det\Lambda\,\overline\psi(x)\gamma^5\psi(x),&&\text{Lorentz pseudoscalar},\\
\overline\psi'(x')\gamma^5\gamma^\mu\psi'(x')
&=\overline\psi(x)S^{-1}\gamma^5SS^{-1}\gamma^\mu S\psi(x)
=\det\Lambda\,\Lambda^\mu{}_{\nu}\overline\psi(x)\gamma^5\gamma^\nu\psi(x),&&\text{Lorentz pseudovector},\\
\overline\psi'(x')\sigma^{\mu\nu}\psi'(x')
&=\overline\psi(x)S^{-1}\sigma^{\mu\nu}S\psi(x)
=\Lambda^\mu{}_{\rho}\Lambda^\nu{}_{\sigma}
\overline\psi(x)\sigma^{\rho\sigma}\psi(x),&&\text{Lorentz tensor}.
\end{align*}
\endgroup
The transformation $S$ may be either continuous or a space-reflection transformation.